\subsection{赋值除环上的赋范空间}

定义 5. 若 E 是非离散赋值除环 K 上的（左）向量空间，E 上的范数是一个映射 \(\mathbf{x} \to p(\mathbf{x})\) ，它把 E 映入 \(\mathbf{R}_+\) ，且满足下列公理：

\[
\left(\mathrm{NO} _ {\mathbf {I I}}\right) \quad p (\boldsymbol {x} + \boldsymbol {y}) \leqslant p (\boldsymbol {x}) + p (\boldsymbol {y}) \text {   for   all   } \boldsymbol {x}, \boldsymbol {y} \text {   in   } \mathbf {E};
\]

\[
\left(\mathrm{NO} _ {\mathbf {I I I}}\right) \quad p (t \boldsymbol {x}) = | t | p (\boldsymbol {x}) \text {for all} t \in \mathrm{K} \text {and all} \boldsymbol {x} \in \mathrm{E}.
\]

最常遇到的赋范空间以 R 或 C 为标量域（带通常绝对值）。

由 \((\mathrm{NO}_{\mathrm{III}})\) 特别推出 \(p(-x)=p(x)\) ；因此若令 \(d(x,y)=p(x-y)\) ，则 d 是加法群 E 上的不变度量，并且定义一个与 E 的加法群结构相容的度量空间拓扑（第 1 号，命题 3）；此外，映射

\[
(t, \boldsymbol {x}) \rightarrow t \boldsymbol {x}
\]

在 \(\mathbf{K} \times \mathbf{E}\) 上连续；因为我们有

\[
t \boldsymbol {x} - t _ {0} \boldsymbol {x} _ {0} = (t - t _ {0}) (\boldsymbol {x} - \boldsymbol {x} _ {0}) + (t - t _ {0}) \boldsymbol {x} + t _ {0} (\boldsymbol {x} - \boldsymbol {x} _ {0})
\]

由此得到

\[
p (t \boldsymbol {x} - t _ {0} \boldsymbol {x} _ {0}) \leqslant | t - t _ {0} | p (\boldsymbol {x} - \boldsymbol {x} _ {0}) + | t - t _ {0} | p (\boldsymbol {x}) + | t _ {0} | p (\boldsymbol {x} - \boldsymbol {x} _ {0}),
\]

这表明只要取 \(|t - t_{0}|\) 与 \(p(x - x_{0})\) 充分小，就可使左端任意小。

定义 6. 若 K 是非离散赋值除环，则 K 上的向量空间 E 赋有由 E 上一个给定范数所定义的结构后，称为 K 上的赋范空间。

赋范空间总被视为赋有由其范数定义的拓扑与一致结构。

例。1) 在非离散赋值除环 K 上，把 K 自身看作（左或右）向量空间，绝对值 \(|x|\) 是一个范数。

\begin{enumerate}
\def\labelenumi{\arabic{enumi})}
\setcounter{enumi}{1}
\item
  表达式 \(||x|| = \sqrt{\sum_{i=1}^{n} x_i^2}\) ——我们曾称之为空间 \(\mathbf{R}^n\) 上的欧氏范数（第 VI 章，§ 2）——显然是定义 5 意义下的范数。函数 \(\sup_{1 \leqslant i \leqslant n} |x_i|\) 与 \(\sum_{i=1}^{n} |x_i|\) 也是范数。
\item
  设 \(\mathcal{B}(\mathrm{E};\mathrm{K})\) 是全体函数 \(f\) （在集 \(\mathrm{E}\) 上、取值于非离散赋值除环 \(\mathbf{K}\) ，且使实值函数 \(x\to |f(x)|\) 在 \(\mathbf{E}\) 上有界）所成的集。这个集显然是（左或右）向量空间 \(\mathbf{K}^{\mathbf{E}}\) （由 \(\mathbf{E}\) 到 \(\mathbf{K}\) 的一切映射组成）的向量子空间。若令 \(p(f)=\sup_{x\in\mathbb{E}}|f(x)|\) ，则 \(p\) 是向量空间 \(\mathcal{B}(\mathrm{E};\mathrm{K})\) 上的范数（参见第 X 章，§ 1）。
\end{enumerate}

* 4) 在向量空间 \(\mathcal{C}(\mathrm{I};\mathbf{R})\) （由定义在区间 \(\mathbf{I} = [0,1]\) （\(\mathbf{R}\) 中的）上的一切有限连续实值函数组成）上，函数

\[
p (x) = \int_ {0} ^ {1} | x (t) | d t
\]

是一个范数。*

在赋范空间 E 中，以 o 为中心、i 为半径的（闭）球 B，即全体 \(x \in E\) 中满足 \(p(x) \leqslant i\) 者所成的集，将称为 E 的单位球。我们来证明：E 中 o 的一个基本邻域系由单位球在位似变换 \(x \to tx\) （t 取遍 K 的非零元素）下的像组成。B 在这位似变换下的像是以 o 为中心、\(|t|\) 为半径的闭球；因此只需证明，对每个实数 r \textgreater{} o，存在 \(t \in K\) 使 \(o < |t| < r\) 。由于 K 的绝对值不是平凡的，存在 \(t_{0} \in K\) 使 \(o < |t_{0}| < i\) ；因此取 \(t = t_{0}^{n}\) （n 为充分大的整数）即可使 \(|t| = |t_{0}|^{n} < r\) 。

定义 7. 向量空间 E（在非离散赋值除环 K 上的）上的两个范数称为等价的，如果它们在 E 上定义相同的拓扑。

命题 7. 两个范数 \(p, q\) （在向量空间 \(E\) 上）等价，当且仅当存在两个数 \(a > 0, b > 0\) 使得

\[
a. p (\pmb {x}) \leqslant q (\pmb {x}) \leqslant b. p (\pmb {x})\tag{I}
\]

对一切 \(x \in E\) 成立。

这些不等式是充分的，因为由关系式 \(a.p(x) \leqslant q(x)\) 推出，对每个 \(r > 0\) ，以 \(o\) 为中心、\(ar\) 为半径的闭球（关于范数 \(q\) ）含于以 \(o\) 为中心、\(r\) 为半径的闭球（关于范数 \(p\) ）；因而由 \(q\) 定义的拓扑比由 \(p\) 定义的拓扑细。类似地，不等式 \(q(x) \leqslant b.p(x)\) 表明由 \(p\) 定义的拓扑比由 \(q\) 定义的拓扑细，故 \(p\) 与 \(q\) 等价。

现在证明不等式 (1) 是必要的。若由 \(q\) 定义的拓扑比由 \(p\) 定义的拓扑细，则关于 \(p\) 的单位球包含一个以 \(o\) 为中心、半径 \(\alpha > 0\) 的闭球（关于 \(q\) 的）；即关系 \(q(x) \leqslant \alpha\) 蕴含 \(p(x) \leqslant 1\) 。若 \(t_0 \in K\) 满足 \(0 < |t_0| < 1\) ，则对每个 \(x \neq 0\) （\(E\) 中的）存在唯一的有理整数 \(k\) 使 \(\alpha |t_0| < q(t_0^k x) \leqslant \alpha\) ；因此 \(p(t_0^k x) \leqslant 1\) ，从而

\[
p (\boldsymbol {x}) \leqslant \frac {\mathrm{I}}{\left| t _ {0} \right| ^ {k}} \leqslant \frac {\mathrm{I}}{\alpha \left| t _ {0} \right|} q (\boldsymbol {x});
\]

令 \(a = \alpha |t_0|\) ，便有 \(a.p(x) \leqslant q(x)\) 对一切 \(x \neq 0\) 成立，且当 \(x = 0\) 时此不等式也成立。类似地可证，若由 \(p\) 定义的拓扑比由 \(q\) 定义的拓扑细，则存在 \(b > 0\) 使 \(q(x) \leqslant b.p(x)\) 。

例。在空间 \(R^{n}\) 中，三个范数

\[
\sqrt {\sum_ {i = 1} ^ {n} x _ {i} ^ {2}}, \quad \sup _ {1 \leqslant i \leqslant n} | x _ {i} | \quad \text { and } \quad \sum_ {i = 1} ^ {n} | x _ {i} |
\]

等价，因为我们有

\[
\sup _ {1 \leqslant i \leqslant n} | x _ {i} | \leqslant \sqrt {\sum_ {i = 1} ^ {n} x _ {i} ^ {2}} \leqslant \sum_ {i = 1} ^ {n} | x | _ {i} \leqslant n. \sup _ {1 \leqslant i \leqslant n} | x _ {i} |.\tag{2}
\]

命题 8. 设 E 是非离散赋值除环上的赋范空间，p 是 E 上的范数，\(\hat{E}\) 是作为加法群 E 的完备化的加法拓扑群。则函数 \((t, x) \to tx\) 可由连续性扩张到 \(\hat{K} \times \hat{E}\) 上，并在 \(\hat{E}\) 上定义一个 \(\hat{K}\) 上的向量空间结构；范数 p 可由连续性扩张为范数 \(\overline{p}\) （在 \(\hat{E}\) 上），它定义 \(\hat{E}\) 的拓扑。

tx 的连续延拓是关于两个交换群之积到第三个交换群中的连续双线性映射的延拓定理的特例（第 III 章，§ 6，第 5 号，定理 1）；由恒等式延拓原理，有 \(\mathbf{I} \cdot \mathbf{x} = \mathbf{x}\) 与 \(t(ux) = (tu)x\) （对 \(t \in \hat{\mathbf{K}}\) 、 \(u \in \hat{\mathbf{K}}\) 与 \(\mathbf{x} \in \hat{\mathbf{E}}\) ）；因而外律 \((t, \mathbf{x}) \to tx\) 确实在 \(\hat{\mathbf{E}}\) 上定义了 \(\hat{\mathbf{K}}\) 上的向量空间结构。另一方面，不变度量 \(\overline{d}(\mathbf{x}, \mathbf{y}) = \overline{p}(\mathbf{x} - \mathbf{y})\) 扩张为不变度量 \(\overline{d}\) （在 \(\hat{\mathbf{E}}\) 上）（ \(\S 2\) ，第 1 号，命题 1），它定义 \(\hat{\mathbf{E}}\) 的拓扑；若令 \(\overline{p}(\mathbf{x}) = \overline{d}(\mathrm{o}, \mathbf{x})\) ，则 \(\overline{p}\) 是 \(p\) 的连续延拓，且满足公理 \((\mathrm{NO}_{\mathbf{I}})\) 与 \((\mathrm{NO}_{\mathbf{II}})\) ；由于 \(tx\) 在 \(\hat{\mathbf{K}} \times \hat{\mathbf{E}}\) 上连续，\(\overline{p}\) 还满足 \((\mathrm{NO}_{\mathbf{III}})\) （恒等式延拓原理），从而是 \(\hat{\mathbf{E}}\) 上的范数。

当我们需要考虑向量空间 E 上一个确定的赋范空间结构时，除非这个记号可能引起混淆，我们通常以 \(\|x\|\) 表示向量 x 的范数。

